\chapter{2001年大纲卷（理）}
\section{单选题}
\begin{problem}
若 \(\sin \theta \cos \theta > 0\)，则 \(\theta\) 在（\quad）

\choices
    {第一、二象限}
    {第一、三象限}
    {第一、四象限}
    {第二、四象限}
\end{problem}

\begin{answer}
B
\end{answer}

\begin{solution}
当 \(\sin\theta>0\) 且 \(\cos\theta>0\) 时，\(\theta\) 在第一象限；当 \(\sin\theta<0\) 且 \(\cos\theta<0\) 时，\(\theta\) 在第三象限. 故 \(\theta\) 在第一、三象限，选 B.
\end{solution}

\begin{problem}
过点 \(A(1, -1)\)，\(B(-1, 1)\) 且圆心在直线 \(x + y - 2 = 0\) 上的圆的方程是（\quad）

\choices
    {\((x - 3)^2 + (y + 1)^2 = 4\)}
    {\((x + 3)^{2} + (y - 1)^{2} = 4\)}
    {\((x - 1)^{2} + (y - 1)^{2} = 4\)}
    {\((x + 1)^{2} + (y + 1)^{2} = 4\)}
\end{problem}

\begin{answer}
C
\end{answer}

\begin{solution}
设圆心为 \(O_0(h,k)\). 因为圆经过 \(A(1,-1)\)，\(B(-1,1)\)，所以
\[
(h-1)^2+(k+1)^2=(h+1)^2+(k-1)^2,
\]
化简得 \(h=k\). 又因为圆心在直线 \(x+y-2=0\) 上，所以 \(h+k=2\)，从而 \(h=k=1\). 圆的半径满足
\[
r^2=(1-1)^2+(1+1)^2=4,
\]
故圆的方程为 \((x-1)^2+(y-1)^2=4\)，选 C.
\end{solution}

\begin{problem}
设 \(\{a_{n}\}\) 是递增等差数列，前三项的和为 12，前三项的积为 48，则它的首项是（\quad）

\choices
    {1}
    {2}
    {4}
    {6}
\end{problem}

\begin{answer}
B
\end{answer}

\begin{solution}
设等差数列的首项为 \(a\)，公差为 \(d\)，则 \(d>0\)，且前三项为 \(a\)，\(a+d\)，\(a+2d\). 由前三项的和得
\[
3(a+d)=12.
\]
所以 \(a+d=4\). 由前三项的积得
\[
a(a+d)(a+2d)=4a(8-a)=48.
\]
即 \(a^2-8a+12=0\)，所以 \(a=2\) 或 \(a=6\). 由于 \(d=4-a>0\)，只能取 \(a=2\)，选 B.
\end{solution}

\begin{problem}
若定义在区间 \((-1,0)\) 内的函数 \(f(x) = \log_{2a}(x + 1)\) 满足 \(f(x) > 0\)，则 \(a\) 的取值范围是（\quad）

\choices
    {\(\left(0, \frac{1}{2}\right)\)}
    {\(\left(0, \frac{1}{2}\right]\)}
    {\(\left(\frac{1}{2}, +\infty\right)\)}
    {\((0, +\infty)\)}
\end{problem}

\begin{answer}
A
\end{answer}

\begin{solution}
当 \(-1<x<0\) 时，\(0<x+1<1\). 要使 \(f(x)=\log_{2a}(x+1)\) 在整个区间上有定义，必须有 \(2a>0\) 且 \(2a\ne1\). 对于 \(0<t<1\)，不等式 \(\log_{2a}t>0\) 成立当且仅当底数满足 \(0<2a<1\). 因此
\[
0<a<\frac12,
\]
选 A.
\end{solution}

\begin{problem}
极坐标方程 \(\rho = 2\sin \left(\theta + \frac{\pi}{4}\right)\) 的图形是（\quad）

\choices
    {\choicebitmap[width=0.15\paperwidth]{img/2001/syllabus_science/q5_optA.png}}
    {\choicebitmap[width=0.15\paperwidth]{img/2001/syllabus_science/q5_optB.png}}
    {\choicebitmap[width=0.15\paperwidth]{img/2001/syllabus_science/q5_optC.png}}
    {\choicebitmap[width=0.15\paperwidth]{img/2001/syllabus_science/q5_optD.png}}
\end{problem}

\begin{answer}
C
\end{answer}

\begin{solution}
由极坐标关系 \(x=\rho\cos\theta\)，\(y=\rho\sin\theta\)，有
\[
\rho^2=2\rho\sin\left(\theta+\frac\pi4\right)=\sqrt2\rho(\sin\theta+\cos\theta)=\sqrt2(x+y).
\]
所以
\[
x^2+y^2=\sqrt2(x+y),
\qquad
\left(x-\frac1{\sqrt2}\right)^2+\left(y-\frac1{\sqrt2}\right)^2=1.
\]
这是圆心在第一象限、半径为 \(1\)，且经过原点的圆，其图形对应选项 C.
\end{solution}

\begin{problem}
函数 \( y = \cos x + 1 \)（\(-\pi \leqslant x \leqslant 0\)）的反函数是（\quad）

\choices
    {\(y = -\arccos (x - 1)\)（\(0\leqslant x\leqslant 2\)）}
    {\(y = \pi -\arccos (x - 1)\)（\(0\leqslant x\leqslant 2\)）}
    {\(y = \arccos (x - 1)\)（\(0\leqslant x\leqslant 2\)）}
    {\(y = \pi +\arccos (x - 1)\)（\(0\leqslant x\leqslant 2\)）}
\end{problem}

\begin{answer}
A
\end{answer}

\begin{solution}
在区间 \([-\pi,0]\) 上，\(\cos x\) 单调递增，且 \(y=\cos x+1\) 的值域为 \([0,2]\). 由 \(y=\cos x+1\) 得 \(\cos x=y-1\)，而反函数的自变量取值在 \([0,2]\) 内，原函数的函数值取值在 \([-\pi,0]\) 内，因此
\[
x=-\arccos(y-1).
\]
交换自变量与函数值，反函数为
\[
y=-\arccos(x-1),\qquad 0\leqslant x\leqslant2,
\]
选 A.
\end{solution}

\begin{problem}
若椭圆经过原点，且焦点为 \(F_{1}(1,0)\)，\(F_{2}(3,0)\)，则其离心率为（\quad）

\choices
    {\( \frac{3}{4} \)}
    {\( \frac{2}{3} \)}
    {\( \frac{1}{2} \)}
    {\( \frac{1}{4} \)}
\end{problem}

\begin{answer}
C
\end{answer}

\begin{solution}
椭圆的两个焦点为 \(F_1(1,0)\)，\(F_2(3,0)\)，所以椭圆中心为 \((2,0)\)，焦半距为 \(c=1\). 原点在椭圆上，故原点到两个焦点的距离之和为
\[
2a=\sqrt{(0-1)^2}+\sqrt{(0-3)^2}=1+3=4,
\]
从而 \(a=2\). 因此离心率为
\[
e=\frac ca=\frac12,
\]
选 C.
\end{solution}

\begin{problem}
若 \(0 < \alpha < \beta < \frac{\pi}{4}\)，\(\sin \alpha + \cos \alpha = a\)，\(\sin \beta + \cos \beta = b\)，则（\quad）

\choices
    {\(a > b\)}
    {\(a < b\)}
    {\(ab < 1\)}
    {\(ab > 2\)}
\end{problem}

\begin{answer}
B
\end{answer}

\begin{solution}
令 \(h(x)=\sin x+\cos x\)，则
\[
h'(x)=\cos x-\sin x.
\]
当 \(0<x<\frac\pi4\) 时，\(\cos x>\sin x\)，所以 \(h'(x)>0\)，即 \(h(x)\) 在 \(\left(0,\frac\pi4\right)\) 上单调递增. 由 \(\alpha<\beta\) 得
\[
a=h(\alpha)<h(\beta)=b,
\]
选 B.
\end{solution}

\begin{problem}
在正三棱柱 \(ABC - A_{1}B_{1}C_{1}\) 中，若 \(AB = \sqrt{2} BB_{1}\)，则 \(AB_{1}\) 与 \(C_{1}B\) 所成的角的大小为（\quad）

\choices
    {\(60^{\circ}\)}
    {\(90^{\circ}\)}
    {\(45^{\circ}\)}
    {\(120^{\circ}\)}
\end{problem}

\begin{answer}
B
\end{answer}

\begin{solution}
设正三棱柱的底面边长为 \(s\)，高为 \(h\). 由题意 \(s=\sqrt2h\)，即 \(h^2=\frac{s^2}{2}\). 取 \(B\) 为原点，令
\[
A=(s,0,0),\quad C=\left(\frac s2,\frac{\sqrt3s}{2},0\right),\quad B_1=(0,0,h),
\]
则 \(C_1=\left(\frac s2,\frac{\sqrt3s}{2},h\right)\). 取两直线方向向量
\[
\overrightarrow{AB_1}=(-s,0,h),\qquad \overrightarrow{BC_1}=\left(\frac s2,\frac{\sqrt3s}{2},h\right).
\]
它们的数量积为
\[
\overrightarrow{AB_1}\cdot\overrightarrow{BC_1}=-\frac{s^2}{2}+h^2=0.
\]
因此 \(AB_1\perp C_1B\)，所成的角为 \(90^{\circ}\)，选 B.
\end{solution}

\begin{problem}
设 \(f(x)\)，\(g(x)\) 都是单调函数，有如下四个命题中，正确的命题是

\begin{circlelist}
\item 若 \( f(x) \) 单调递增，\( g(x) \) 单调递增，则 \( f(x) - g(x) \) 单调递增；
\item 若 \( f(x) \) 单调递增，\( g(x) \) 单调递减，则 \( f(x)-g(x) \) 单调递增；
\item 若 \( f(x) \) 单调递减，\( g(x) \) 单调递增，则 \( f(x)-g(x) \) 单调递减；
\item 若 \( f(x) \) 单调递减，\( g(x) \) 单调递减，则 \( f(x)-g(x) \) 单调递减.
\end{circlelist}

（\quad）

\choices
    {\circled{1}\circled{3}}
    {\circled{1}\circled{4}}
    {\circled{2}\circled{3}}
    {\circled{2}\circled{4}}
\end{problem}

\begin{answer}
C
\end{answer}

\begin{solution}
设 \(x_1<x_2\)，记 \(h(x)=f(x)-g(x)\). 若 \(f\) 单调递增、\(g\) 单调递减，则
\[
h(x_2)-h(x_1)=\bigl(f(x_2)-f(x_1)\bigr)+\bigl(g(x_1)-g(x_2)\bigr)>0,
\]
所以命题②正确. 若 \(f\) 单调递减、\(g\) 单调递增，则
\[
h(x_2)-h(x_1)=\bigl(f(x_2)-f(x_1)\bigr)-\bigl(g(x_2)-g(x_1)\bigr)<0,
\]
所以命题③正确.

命题①和命题④均不一定成立. 例如，取 \(f(x)=x\)，\(g(x)=2x\)，则二者都递增而 \(f(x)-g(x)=-x\) 递减；取 \(f(x)=-x\)，\(g(x)=-2x\)，则二者都递减而 \(f(x)-g(x)=x\) 递增. 因此正确的命题为②③，选 C.
\end{solution}

\begin{problem}
一间民房的屋顶有如图三种不同的盖法，分别是：

\begin{circlelist}
\item 单向倾斜；
\item 双向倾斜；
\item 四向倾斜.
\end{circlelist}

记三种盖法屋顶面积分别为 \(P_{1}\)，\(P_{2}\)，\(P_{3}\). 若屋顶斜面与水平面所成的角都是 \(\alpha\)，则（\quad）

\begin{center}
\bitmapfigure[max width=0.70\linewidth]{img/2001/syllabus_science/q11_fig1.png}
\end{center}

\choices
    {\(P_{3} > P_{2} > P_{1}\)}
    {\(P_{3} > P_{2} = P_{1}\)}
    {\(P_{3} = P_{2} > P_{1}\)}
    {\(P_{3} = P_{2} = P_{1}\)}
\end{problem}

\begin{answer}
D
\end{answer}

\begin{solution}
设三种屋顶的水平投影都是同一个屋顶平面，记其面积为 \(S\). 每个斜面与水平面所成的角都为 \(\alpha\)，所以每个斜面的水平投影面积等于该斜面面积乘以 \(\cos\alpha\). 将同一种盖法的各个斜面面积相加，均有
\[
P_i\cos\alpha=S\qquad(i=1,2,3).
\]
由于 \(\cos\alpha>0\)，得到
\[
P_1=P_2=P_3.
\]
故选 D.
\end{solution}

\begin{problem}
如图，小圆圈表示网络的结点，结点之间的连线表示它们有网线相联. 连线标注的数字表示该段网线单位时间内可以通过的最大信息量. 现从结点 \(A\) 向结点 \(B\) 传递信息，信息可以分开沿不同的路线同时传递. 则单位时间内传递的最大信息量为（\quad）

\bitmapinlinefigure[6.4cm][max width=6.4cm]{img/2001/syllabus_science/q12_fig1.png}

\choices
    {26}
    {24}
    {20}
    {19}
\end{problem}

\begin{answer}
D
\end{answer}

\begin{solution}
从 \(B\) 到 \(A\) 的上方两条支路分别受到容量 \(3\) 和 \(4\) 的边限制，合计至多传递 \(3+4=7\). 这两条支路汇合后通向 \(A\) 的边容量为 \(12\)，因此不会限制这 \(7\) 个单位的信息. 中间支路的容量受标为 \(7\) 和 \(6\) 的两条边共同限制，至多传递 \(6\)；下方支路的容量受标为 \(6\) 的边限制，至多传递 \(6\). 所以总流量至多为
\[
7+6+6=19.
\]
沿上方两条支路分别传递 \(3\) 和 \(4\)，沿中间支路传递 \(6\)，沿下方支路传递 \(6\)，各边容量均不超限，确实可以达到 \(19\). 故最大信息量为 \(19\)，选 D.
\end{solution}

\section{填空题}
\begin{problem}
若一个圆锥的轴截面是等边三角形，其面积为 \(\sqrt{3}\)，则这个圆锥的侧面积是\fillinblank{}.
\end{problem}

\begin{answer}
\(2\pi\)
\end{answer}

\begin{solution}
设圆锥的底面半径为 \(r\)，母线长为 \(l\). 轴截面为等边三角形时，轴截面的底边为圆锥底面的直径，所以 \(l=2r\). 轴截面面积为
\[
\frac{\sqrt3}{4}(2r)^2=\sqrt3r^2=\sqrt3,
\]
故 \(r=1\)，从而 \(l=2\). 圆锥的侧面积为
\[
\pi rl=\pi\cdot1\cdot2=2\pi.
\]
\end{solution}

\begin{problem}
双曲线 \(\frac{x^2}{9} - \frac{y^2}{16} = 1\) 的两个焦点为 \(F_{1}\)，\(F_{2}\)，点 \(P\) 在双曲线上，若 \(PF_{1} \perp PF_{2}\)，则点 \(P\) 到 \(x\) 轴的距离为\fillinblank{}.
\end{problem}

\begin{answer}
\(\frac{16}{5}\)
\end{answer}

\begin{solution}
双曲线的参数为 \(a=3\)，\(b=4\)，所以 \(c=\sqrt{a^2+b^2}=5\)，焦点为 \((-5,0)\) 和 \((5,0)\). 设 \(P(x,y)\)，则
\[
\overrightarrow{PF_1}=(x+5,y),\qquad \overrightarrow{PF_2}=(x-5,y).
\]
由 \(PF_1\perp PF_2\)，得
\[
(x+5)(x-5)+y^2=0,qquad x^2+y^2=25.
\]
另一方面，双曲线方程可写为 \(16x^2-9y^2=144\). 代入 \(x^2=25-y^2\)，得
\[
16(25-y^2)-9y^2=144,qquad 25y^2=256.
\]
因此点 \(P\) 到 \(x\) 轴的距离为 \(|y|=\frac{16}{5}\).
\end{solution}

\begin{problem}
设 \( \{a_{n}\} \) 是公比为 \(q\) 的等比数列，\( S_{n} \) 是它的前 \(n\) 项和. 若 \( \{S_{n}\} \) 是等差数列，则 \(q = \)\fillinblank{}.
\end{problem}

\begin{answer}
1
\end{answer}

\begin{solution}
设等比数列首项为 \(a_1\)，公比为 \(q\). 由
\[
S_1=a_1,qquad S_2=a_1(1+q),qquad S_3=a_1(1+q+q^2)
\]
以及 \(\{S_n\}\) 为等差数列，知 \(S_1+S_3=2S_2\). 因此
\[
a_1\bigl[1+(1+q+q^2)-2(1+q)\bigr]=0,
\]
即 \(a_1q(q-1)=0\). 等比数列的首项及公比满足定义中的非零条件，所以 \(a_1\ne0\)，\(q\ne0\)，从而 \(q=1\).
\end{solution}

\begin{problem}
圆周上有 \(2n\) 个等分点（\(n > 1\)），以其中三个点为顶点的直角三角形的个数为\fillinblank{}.
\end{problem}

\begin{answer}
\(2n(n-1)\)
\end{answer}

\begin{solution}
圆周上三个等分点所成的三角形为直角三角形，当且仅当其斜边是该圆的直径. \(2n\) 个等分点中共有 \(n\) 对互为对径点，即有 \(n\) 条直径. 每确定一条直径作为斜边，还可从其余 \(2n-2\) 个等分点中任取一个作为第三个顶点，并且每个直角三角形只有一条直径作为斜边. 因此直角三角形的个数为
\[
n(2n-2)=2n(n-1).
\]
\end{solution}

\section{解答题}
\begin{problem}
如图，在底面是直角梯形的四棱锥 \(S - ABCD\) 中，\(\angle ABC = 90^{\circ}\)，\(SA \perp\) 面 \(ABCD\)，\(SA = AB = BC = 1\)，\(AD = \frac{1}{2}\).

\begin{enumerate}
\item 求四棱锥 \(S - ABCD\) 的体积；
\item 求面 \(SCD\) 与面 \(SBA\) 所成的二面角的正切值.
\end{enumerate}

\bitmapfigure[max width=6.4cm]{img/2001/syllabus_science/q17_fig1.png}
\end{problem}

\begin{solution}
\begin{enumerate}
\item 因为底面是直角梯形，且 \(AD\parallel BC\)，所以 \(AB\) 是梯形的高，从而底面面积为
\[
S_{ABCD}=\frac{AD+BC}{2}\cdot AB=\frac{\frac12+1}{2}\cdot1=\frac34
\]
又因为 \(SA\perp\) 面 \(ABCD\)，所以四棱锥的高为 \(SA=1\)，因此
\[
V_{S-ABCD}=\frac13S_{ABCD}\cdot SA=\frac13\cdot\frac34\cdot1=\frac14
\]
\item 建立空间直角坐标系，取 \(A(0,0,0)\)，\(B(0,1,0)\)，\(C(1,1,0)\)，\(D\left(\frac12,0,0\right)\)，\(S(0,0,1)\). 则平面 \(SBA\) 的一个法向量为 \(\bs{n}_1=(1,0,0)\)，又 \(\overrightarrow{SC}=(1,1,-1)\)，\(\overrightarrow{SD}=\left(\frac12,0,-1\right)\)，所以平面 \(SCD\) 的一个法向量可取
\[
\bs{n}_2=2\left(\overrightarrow{SC}\times\overrightarrow{SD}\right)=(2,-1,1)
\]
设两平面所成的锐二面角为 \(\varphi\)，则
\[
\cos\varphi=\frac{|\bs{n}_1\cdot\bs{n}_2|}{|\bs{n}_1||\bs{n}_2|}=\frac{2}{\sqrt6}=\sqrt{\frac23}
\]
于是
\[
\tan\varphi=\frac{\sqrt{1-\cos^2\varphi}}{\cos\varphi}=\frac{1}{\sqrt2}
\]
\end{enumerate}
\end{solution}

\begin{problem}
已知复数 \(z_{1} = \i(1 - \i)^{3}\).

\begin{enumerate}
\item 求 \(\arg z_{1}\) 及 \(|z_{1}|\)；
\item 当复数 \(z\) 满足 \(|z| = 1\)，求 \(|z - z_1|\) 的最大值.
\end{enumerate}
\end{problem}

\begin{solution}
\begin{enumerate}
\item 先计算得 \((1-\i)^2=-2\i\)，从而 \((1-\i)^3=(-2\i)(1-\i)=-2-2\i\). 所以
\[
z_1=\i(-2-2\i)=2-2\i=2\sqrt2\left(\cos\left(-\frac\pi4\right)+\i\sin\left(-\frac\pi4\right)\right)
\]
因此 \(\arg z_1=-\frac\pi4+2k\pi\)（\(k\in\Z\)），主值为 \(-\frac\pi4\)，且 \(|z_1|=2\sqrt2\).
\item 由三角不等式，\(|z-z_1|\leqslant |z|+|z_1|=1+2\sqrt2\). 当 \(z=-\dfrac{z_1}{|z_1|}=\dfrac{-1+\i}{\sqrt2}\) 时，\(|z|=1\)，且 \(z\) 与 \(-z_1\) 同向，三角不等式取等号. 所以 \(|z-z_1|\) 的最大值为 \(1+2\sqrt2\).
\end{enumerate}
\end{solution}

\begin{problem}
设抛物线 \( y^{2} = 2px \)（\( p > 0 \)）的焦点为 \(F\)，经过点 \(F\) 的直线交抛物线于 \(A\)，\(B\) 两点，点 \(C\) 在抛物线的准线上，且 \( BC \parallel x \) 轴. 证明直线 \(AC\) 经过原点 \(O\).
\end{problem}

\begin{solution}
令 \(u\)，\(v\) 分别为点 \(A\)，\(B\) 在抛物线参数表示中的参数. 令 \(a=\frac p2\)，则抛物线可表示为
\[
(x,y)=(at^2,2at)=\left(\frac p2t^2,pt\right)
\]
其焦点为 \(F(a,0)\)，准线为 \(x=-a\)，因而 \(A=(au^2,2au)\)，\(B=(av^2,2av)\).
由于 \(A\)，\(B\)，\(F\) 共线，有
\[
\begin{vmatrix}
au^2-a & 2au\\
av^2-a & 2av
\end{vmatrix}=0
\]
化简得 \(a^2(u-v)(uv+1)=0\). 因为 \(a>0\) 且 \(A\ne B\)，所以 \(uv=-1\).

由 \(BC\parallel x\) 轴，点 \(C\) 与点 \(B\) 的纵坐标相同；又 \(C\) 在准线上，所以 \(C=(-a,2av)\). 于是
\[
\begin{vmatrix}
au^2 & 2au\\
-a & 2av
\end{vmatrix}=2a^2u(uv+1)=0
\]
这说明向量 \(\overrightarrow{OA}\) 与 \(\overrightarrow{OC}\) 共线，故直线 \(AC\) 经过原点 \(O\).
\end{solution}

\begin{problem}
已知 \(i\)，\(m\)，\(n\) 是正整数，且 \(1 < i \leqslant m < n\).

\begin{enumerate}
\item 证明： \(n^i \mathrm A_{m}^{i} < m^i \mathrm A_{n}^{i}\)；
\item 证明： \((1+m)^{n}>(1+n)^{m}\).
\end{enumerate}
\end{problem}

\begin{solution}
\begin{enumerate}
\item 按排列数的定义，\(\displaystyle\frac{\mathrm A_m^i}{m^i}=\prod_{k=0}^{i-1}\frac{m-k}{m}=\prod_{k=1}^{i-1}\left(1-\frac{k}{m}\right)\)，同理 \(\displaystyle\frac{\mathrm A_n^i}{n^i}=\prod_{k=0}^{i-1}\frac{n-k}{n}=\prod_{k=1}^{i-1}\left(1-\frac{k}{n}\right)\). 对每个 \(k=1,2,\ldots,i-1\)，由 \(m<n\) 得 \(1-\frac{k}{m}<1-\frac{k}{n}\). 又因为 \(i\leqslant m<n\)，上述各因子均为正；并且 \(i>1\) 保证至少有一个严格不等号，所以
\[
\frac{\mathrm A_m^i}{m^i}<\frac{\mathrm A_n^i}{n^i}\Longrightarrow n^i\mathrm A_m^i<m^i\mathrm A_n^i
\]
\item 设 \(h(x)=\dfrac{\ln(1+x)}{x}\)（\(x>0\)），则 \(h'(x)=\frac{\frac{x}{1+x}-\ln(1+x)}{x^2}\). 因为
\[
\ln(1+x)=\int_0^x\frac{1}{1+t}\,\mathrm dt>\frac{x}{1+x}
\]
所以 \(h'(x)<0\)，即 \(h(x)\) 在 \((0,+\infty)\) 上单调递减. 由 \(m<n\) 得 \(\frac{\ln(1+m)}{m}>\frac{\ln(1+n)}{n}\). 两边同乘 \(mn>0\)，再利用指数函数的单调性，得到
\[
n\ln(1+m)>m\ln(1+n)\Longrightarrow (1+m)^n>(1+n)^m
\]
\end{enumerate}
\end{solution}

\begin{problem}
从社会效益和经济效益出发，某地投入资金进行生态环境建设，并以此发展旅游产业. 根据规划，本年度投入 \(800\text{ 万元}\)，以后每年投入将比上年减少 \(\frac{1}{5}\). 本年度当地旅游业收入估计为 \(400\text{ 万元}\)，由于该项建设对旅游业的促进作用，预计今后的旅游业收入每年会比上年增加 \(\frac{1}{4}\).

\begin{enumerate}
\item 设 \(n\) 年内（本年度为第一年）总投入为 \(a_{n}\text{ 万元}\)，旅游业总收入为 \(b_{n}\text{ 万元}\). 写出 \(a_{n}\)，\(b_{n}\) 的表达式；
\item 至少经过几年旅游业的总收入才能超过总投入？
\end{enumerate}
\end{problem}

\begin{solution}
\begin{enumerate}
\item 每年投入构成首项为 \(800\)，公比为 \(\frac45\) 的等比数列，所以
\[
a_n=800\frac{1-\left(\frac45\right)^n}{1-\frac45}=4000\left[1-\left(\frac45\right)^n\right]
\]
每年旅游业收入构成首项为 \(400\)，公比为 \(\frac54\) 的等比数列，所以
\[
b_n=400\frac{\left(\frac54\right)^n-1}{\frac54-1}=1600\left[\left(\frac54\right)^n-1\right]
\]
\item 令 \(t=\left(\frac54\right)^n>1\)，则 \(\left(\frac45\right)^n=\frac1t\). 条件 \(b_n>a_n\) 等价于
\[
1600(t-1)>4000\left(1-\frac1t\right)
\Longleftrightarrow (t-1)(4t-10)>0
\]
由于 \(t>1\)，所以必须有 \(t>\frac52\)，即
\[
\left(\frac54\right)^n>\frac52
\]
又 \(\left(\frac54\right)^4=\frac{625}{256}<\frac52\)，而
\(\left(\frac54\right)^5=\frac{3125}{1024}>\frac52\)，
故至少经过 \(5\) 年旅游业的总收入才能超过总投入.
\end{enumerate}
\end{solution}

\begin{problem}
设 \( f(x) \) 是定义在 \( \R \) 上的偶函数，其图象关于直线 \( x = 1 \) 对称，对任意 \( x_{1} \)，\( x_{2} \in \left[0, \frac{1}{2}\right] \)，都有 \( f(x_{1} + x_{2}) = f(x_{1}) \cdot f(x_{2}) \)，且 \( f(1) = a > 0 \).

\begin{enumerate}
\item 求 \(f\left(\frac{1}{2}\right)\)，\(f\left(\frac{1}{4}\right)\)；
\item 证明： \( f(x) \) 是周期函数；
\item 记 \(a_{n} = f\left(2n + \frac{1}{2n}\right)\)，求 \(\displaystyle\lim_{n \to \infty}(\ln a_{n})\).
\end{enumerate}
\end{problem}

\begin{solution}
\begin{enumerate}
\item 取 \(x_1=x_2=\frac12\)，由已知得 \(f(1)=f\left(\frac12\right)^2=a\). 对任意 \(x\in\left[0,\frac12\right]\)，取 \(x_1=x_2=\frac{x}{2}\)，则 \(f(x)=f\left(\frac{x}{2}\right)^2\geqslant0\)，因此 \(f\left(\frac12\right)=\sqrt a\). 再取 \(x_1=x_2=\frac14\)，得 \(f\left(\frac12\right)=f\left(\frac14\right)^2\). 结合 \(f\left(\frac14\right)\geqslant0\)，可得 \(f\left(\frac14\right)=a^{1/4}\). 又由 \(f\left(\frac12\right)=f(0)f\left(\frac12\right)\) 以及 \(f\left(\frac12\right)>0\)，得 \(f(0)=1\).
\item 由图象关于直线 \(x=1\) 对称，\(f(x)=f(2-x)\). 又因为 \(f\) 是偶函数，\(f(2-x)=f(-(2-x))=f(x-2)\). 所以对任意 \(x\in\R\)，都有 \(f(x)=f(x-2)\)，即 \(f(x)\) 是以 \(2\) 为一个周期的周期函数.
\item 由周期性，\(a_n=f\left(2n+\frac{1}{2n}\right)=f\left(\frac{1}{2n}\right)\). 反复使用已知的乘法关系，在 \(k=1,2,\ldots,n\) 时均有 \(f\left(\frac{k}{2n}\right)=f\left(\frac{k-1}{2n}\right)f\left(\frac{1}{2n}\right)\)，其中两个自变量都属于 \(\left[0,\frac12\right]\). 因而 \(f\left(\frac12\right)=f\left(\frac{1}{2n}\right)^n\). 由第（1）问，左边为 \(\sqrt a>0\)，且 \(f\left(\frac{1}{2n}\right)\geqslant0\)，所以 \(f\left(\frac{1}{2n}\right)=a^{1/(2n)}\). 于是
\[
\ln a_n=\frac{\ln a}{2n}\longrightarrow0
\]
\end{enumerate}
\end{solution}
